% Options for packages loaded elsewhere
\PassOptionsToPackage{unicode}{hyperref}
\PassOptionsToPackage{hyphens}{url}
\documentclass[
  journal]{IEEEtran}
\usepackage{xcolor}
\usepackage{amsmath,amssymb}
\setcounter{secnumdepth}{-\maxdimen} % remove section numbering
\usepackage{iftex}
\ifPDFTeX
  \usepackage[T1]{fontenc}
  \usepackage[utf8]{inputenc}
  \usepackage{textcomp} % provide euro and other symbols
\else % if luatex or xetex
  \usepackage{unicode-math} % this also loads fontspec
  \defaultfontfeatures{Scale=MatchLowercase}
  \defaultfontfeatures[\rmfamily]{Ligatures=TeX,Scale=1}
\fi
\usepackage{lmodern}
\ifPDFTeX\else
  % xetex/luatex font selection
\fi
% Use upquote if available, for straight quotes in verbatim environments
\IfFileExists{upquote.sty}{\usepackage{upquote}}{}
\IfFileExists{microtype.sty}{% use microtype if available
  \usepackage[]{microtype}
  \UseMicrotypeSet[protrusion]{basicmath} % disable protrusion for tt fonts
}{}
\makeatletter
\@ifundefined{KOMAClassName}{% if non-KOMA class
  \IfFileExists{parskip.sty}{%
    \usepackage{parskip}
  }{% else
    \setlength{\parindent}{0pt}
    \setlength{\parskip}{6pt plus 2pt minus 1pt}}
}{% if KOMA class
  \KOMAoptions{parskip=half}}
\makeatother
\usepackage{longtable,booktabs,array}
\usepackage{tabularx}
\setlength{\tabcolsep}{2.5pt}
\renewcommand{\arraystretch}{1.05}
\newcounter{none} % for unnumbered tables
\usepackage{calc} % for calculating minipage widths
% Correct order of tables after \paragraph or \subparagraph
\usepackage{etoolbox}
\makeatletter
\patchcmd\longtable{\par}{\if@noskipsec\mbox{}\fi\par}{}{}
\makeatother
% Allow footnotes in longtable head/foot
\IfFileExists{footnotehyper.sty}{\usepackage{footnotehyper}}{\usepackage{footnote}}
\makesavenoteenv{longtable}
\setlength{\emergencystretch}{3em} % prevent overfull lines
\providecommand{\tightlist}{%
  \setlength{\itemsep}{0pt}\setlength{\parskip}{0pt}}
\usepackage{bookmark}
\IfFileExists{xurl.sty}{\usepackage{xurl}}{} % add URL line breaks if available
\urlstyle{same}
% fallback for those not using the hyperref driver hyperxmp:
\makeatletter
\@ifundefined{xmpquote}{\newcommand{\xmpquote}[1]{#1}}{}
\makeatother
\hypersetup{
  hidelinks,
  pdfcreator={LaTeX via pandoc}}

\author{}
\date{}

\begin{document}

\section{Autonomous UAV Navigation and Localization in GPS/GNSS-Denied
Environments:}\label{autonomous-uav-navigation-and-localization-in-gpsgnss-denied-environments}

\section{A Systematic Literature Review of Sensor-Fusion
Approaches}\label{a-systematic-literature-review-of-sensor-fusion-approaches}

\subsection{Version-1 Full-Text
Synthesis}\label{version-1-full-text-synthesis}

\textbf{Review window:} 2010--2026\\
\textbf{Screened-included set:} 636 studies\\
\textbf{Version-1 full-text corpus:} 171 studies (26.9\%)\\
\textbf{Analysis date:} September 2026

\subsection{Abstract}\label{abstract}

GPS/GNSS denial caused by indoor occlusion, urban canyons, underground
operation, foliage, jamming, or spoofing creates a fundamental
localization problem for autonomous unmanned aerial vehicles (UAVs).
This systematic literature review examines sensor-fusion approaches for
navigation and localization under such conditions. The protocol defines
a 2010--2026 review window, a UAV-focused population, multi-sensor
interventions, localization and robustness outcomes, and indoor, urban,
underground, forest, maritime, and electronic-warfare contexts. The
review workflow screened 636 studies as included and completed
machine-assisted full-text extraction for 171 locally available papers.
The Version-1 full-text corpus is therefore an intermediate subset, not
the final complete review population. Descriptively, the extracted
studies are concentrated in visual-inertial odometry,
visual-LiDAR-inertial fusion, LiDAR-inertial odometry/SLAM, and
radio-based positioning. Evaluation modality counts are 22
real-world-only, 78 simulation-only, and 71 both real-world and
simulation studies. Quality appraisal assigns 38 studies to Q-high, 84
to Q-medium, and 49 to Q-low using a 10-point rubric. The evidence
supports a strong role for complementary sensing but does not support a
pooled accuracy ranking because metrics, units, datasets, and reporting
completeness are heterogeneous and many numeric values are not reported
in the structured table. The main conclusions are consequently
descriptive and should be updated when additional full texts become
available.

\textbf{Keywords:} UAV navigation; GNSS-denied; GPS-denied;
visual-inertial odometry; LiDAR-inertial odometry; SLAM; sensor fusion;
systematic review.

\subsection{1. Introduction}\label{1-introduction}

Reliable position estimation is a prerequisite for autonomous UAV
flight. When GPS or another GNSS service is unavailable, degraded,
jammed, spoofed, or unreliable because of indoor or built-up
environments, a vehicle must construct its state estimate from onboard
sensing and prior information. Failure can produce drift, loss of scale
or observability, collision risk, mission failure, or unsafe return
behavior. The problem is particularly difficult in texture-poor
interiors, urban canyons, tunnels, mines, forests, smoke or dust, and
adversarial electronic environments.

Sensor fusion is central because no single sensing modality is reliable
in all of these conditions. Cameras provide geometric and semantic
information but are sensitive to illumination, texture, motion blur, and
occlusion. Inertial measurements provide high-rate short-term motion
information but drift when not corrected. LiDAR can provide range
geometry in low-light or texture-poor scenes, although it adds payload,
power, and computational requirements. Radio, UWB, radar, sonar, and
map-based information can add complementary constraints but may require
infrastructure, environmental assumptions, or additional calibration.

The review protocol addresses a gap in systematic comparison of these
configurations for UAV navigation in GPS/GNSS-denied environments. It
asks:

\textbf{RQ1.} What sensor-fusion configurations are most commonly
used?\\
\textbf{RQ2.} What performance metrics and accuracy levels are achieved
in different environments?\\
\textbf{RQ3.} What algorithmic approaches and comparative advantages are
reported?\\
\textbf{RQ4.} What limitations and future research directions remain?

This manuscript reports a \textbf{Version-1 full-text synthesis}. The
canonical screened-included set contains 636 papers, but only 171
corresponding PDFs were locally available for structured full-text
extraction at the time of analysis. The conclusions therefore
characterize the available full-text subset and must not be generalized
without qualification to all 636 included records.

\subsection{2. Methods}\label{2-methods}

\subsubsection{2.1 Review design and
protocol}\label{21-review-design-and-protocol}

The review follows the repository research protocol
(\texttt{00\_scope/research\_protocol.md}). The population is UAVs,
drones, and closely related aerial autonomous systems. The intervention
is multi-sensor fusion, including combinations of vision, LiDAR, IMU,
radar, UWB, and related aiding sensors. Comparisons include baselines
and, where reported, GPS-assisted or alternative navigation conditions.
Outcomes include localization accuracy, trajectory error, drift,
robustness, computational efficiency, and mission success. The context
includes indoor, urban canyon, underground, forest, maritime, jamming,
and spoofing environments.

\subsubsection{2.2 Search and
eligibility}\label{22-search-and-eligibility}

The protocol specifies IEEE Xplore, Scopus, Web of Science, and
supplementary Google Scholar searching. Its principal search concepts
combine GPS/GNSS denial or indoor/underground navigation with
localization, navigation, or SLAM and UAV, drone, or robot terms. The
publication window is 2010--2026.

Included studies are English-language, peer-reviewed journal or
conference papers that address GPS/GNSS-denied navigation, focus on
UAVs/drones/robots, present a sensor-fusion approach, and provide
experimental validation or simulation results. The protocol excludes
purely theoretical work without validation, GPS-only navigation,
terrestrial-only vehicles, single-sensor-only studies, non-English
publications, and work published before 2010.

\subsubsection{2.3 Selection and Version-1 corpus
construction}\label{23-selection-and-version-1-corpus-construction}

The repository records 636 screened-included studies. The full-text
corpus was constructed by matching locally available PDF stems in
\texttt{05\_papers\_fulltext/} to study IDs in
\texttt{screened\_included\_v2.csv}. This produced 171 matched records,
equivalent to 26.9\% of the screened-included set. The denominator for
full-text coverage is therefore the 636 included studies.

\subsubsection{2.4 Data extraction}\label{24-data-extraction}

The authoritative table is
\texttt{02\_data\_processed/extracted\_master\_v2.csv}. It contains
study identity, authors, year, DOI, venue, full-text status, page count,
platform type, sensors, primary method, method category, environment,
experiment type, real/simulation status, reported metrics, ATE/RMSE
field, application domain, multi-agent status, notes, four quality
components, total QA score, QA tier, and citation tier.

Extraction was machine-assisted from PDF text using the repository
Phase-7 runner. The resulting labels are suitable for a reproducible
descriptive baseline but are not equivalent to independent human coding.
A 20-paper stratified validation sample was generated in
\texttt{08\_docs/extraction\_validation\_sample\_v1.csv}; its
human-review fields remain to be completed unless subsequently updated.

\subsubsection{2.5 Quality appraisal}\label{25-quality-appraisal}

Each extracted paper receives four components:

{\def\LTcaptype{none} % do not increment counter
\begin{tabularx}{\columnwidth}{@{}XrX@{}}
\toprule\noalign{}
Component & Range & Interpretation \\
\midrule\noalign{}
\endhead
\bottomrule\noalign{}
\endlastfoot
Experimental rigor & 0--4 & Design clarity, validation, ground truth,
and repeatability \\
Reporting completeness & 0--3 & Metrics, scale, hardware/data, and
failure reporting \\
Baseline fairness & 0--2 & Baseline presence and matched comparison \\
Reproducibility & 0--1 & Code, data, or sufficiently detailed
parameters \\
\textbf{Total} & \textbf{0--10} & Sum of the four components \\
\end{tabularx}
}

The tiers are Q-high (8--10), Q-medium (5--7), and Q-low (0--4). A
quality tier is a reporting and evidence-quality indicator; it is not
proof that the proposed method is superior.

\subsubsection{2.6 Synthesis}\label{26-synthesis}

The synthesis is descriptive. Frequencies and cross-tabulations are
reported for methods, sensors, environments, applications, experiment
modality, and QA tier. Environment and sensor fields are multi-valued,
so their frequency sums can exceed 171. Because the extracted
performance values use heterogeneous metrics, units, datasets, and
evaluation protocols, no pooled meta-analysis or universal method
ranking is attempted.

\subsection{3. Results}\label{3-results}

\subsubsection{3.1 Selection and corpus
profile}\label{31-selection-and-corpus-profile}

{\def\LTcaptype{none} % do not increment counter
\begin{tabularx}{\columnwidth}{@{}Xr@{}}
\toprule\noalign{}
Stage & Records \\
\midrule\noalign{}
\endhead
\bottomrule\noalign{}
\endlastfoot
Screened-included studies & 636 \\
Full-text PDFs locally available and matched & 171 \\
Full-text coverage of included set & 26.9\% \\
\end{tabularx}
}

The Version-1 corpus is thus a convenience-of-access subset of the
included set. It is appropriate for an interim synthesis and pipeline
validation, but not for final prevalence estimates over all included
studies.

\subsubsection{3.2 Publication years and
platforms}\label{32-publication-years-and-platforms}

The 171 extracted records span 2013--2026. The largest annual group in
the current corpus is 2026 (36 records), followed by 2018 (24), 2019
(21), and 2022 and 2023 (19 each). These counts should be interpreted as
the temporal distribution of the currently available PDFs, not as the
publication trend of the full 636-paper set.

The platform field identifies 95 records as UAV and 76 as
multi-platform. Because the latter label can represent papers whose text
mentions or evaluates more than one platform, it should not be read as
evidence that 76 studies are non-UAV studies.

\subsubsection{3.3 Sensor
configurations}\label{33-sensor-configurations}

{\def\LTcaptype{none} % do not increment counter
\begin{longtable}[]{@{}lr@{}}
\toprule\noalign{}
Sensor field & Papers \\
\midrule\noalign{}
\endhead
\bottomrule\noalign{}
\endlastfoot
IMU & 170 \\
3-D LiDAR & 93 \\
Stereo camera & 87 \\
mmWave radar & 48 \\
Monocular camera & 46 \\
UWB & 46 \\
Sonar & 34 \\
Optical flow & 30 \\
2-D LiDAR & 25 \\
Thermal camera & 23 \\
Depth camera & 6 \\
\end{longtable}
}

These are multi-label counts. IMU appears in nearly all extracted
records, while cameras and LiDAR are frequent complementary sources. The
presence of a sensor in a paper does not establish that it was the
primary source of accuracy or that it was evaluated under the same
conditions as another paper.

\subsubsection{3.4 Algorithmic taxonomy}\label{34-algorithmic-taxonomy}

{\def\LTcaptype{none} % do not increment counter
\begin{longtable}[]{@{}lr@{}}
\toprule\noalign{}
Primary method & Papers \\
\midrule\noalign{}
\endhead
\bottomrule\noalign{}
\endlastfoot
Visual-LiDAR-inertial fusion & 73 \\
Visual-inertial odometry & 61 \\
LiDAR-inertial odometry/SLAM & 19 \\
Radio-based positioning & 9 \\
Thermal-inertial odometry & 4 \\
Visual SLAM & 3 \\
Hybrid inertial fusion & 2 \\
\end{longtable}
}

The largest categories combine complementary geometric and inertial
information. Visual-inertial methods offer compact, high-rate estimation
but remain vulnerable to visual degradation and accumulated drift.
LiDAR-inertial methods provide direct geometric constraints and can be
advantageous in texture-poor or low-light settings, at the cost of
payload and processing requirements. Radio-based methods can provide
global or infrastructure-linked constraints, but their effectiveness
depends on anchor placement, propagation, or environmental availability.
These are trade-offs represented in the current records, not a universal
performance ordering.

\subsubsection{3.5 Environments and
applications}\label{35-environments-and-applications}

{\def\LTcaptype{none} % do not increment counter
\begin{longtable}[]{@{}lr@{}}
\toprule\noalign{}
Environment label & Papers \\
\midrule\noalign{}
\endhead
\bottomrule\noalign{}
\endlastfoot
Adversarial/EW & 168 \\
Urban & 154 \\
Indoor & 144 \\
Underground/tunnel & 136 \\
Forest/canopy & 89 \\
Maritime/water & 66 \\
Mixed/unstructured & 1 \\
\end{longtable}
}

The unusually high multi-label counts show why environment categories
must not be added as mutually exclusive classes. They reflect the labels
produced by the current extraction process and require human validation
before being used for precise prevalence claims.

Application labels are inspection (107), mapping/surveying (49),
search-and-rescue (11), general navigation (3), and defense/tactical
(1). Inspection and mapping are therefore the most frequent application
labels in the Version-1 table.

\subsubsection{3.6 Evaluation modality and reported
metrics}\label{36-evaluation-modality-and-reported-metrics}

{\def\LTcaptype{none} % do not increment counter
\begin{longtable}[]{@{}lr@{}}
\toprule\noalign{}
Evaluation modality & Papers \\
\midrule\noalign{}
\endhead
\bottomrule\noalign{}
\endlastfoot
Simulation only & 78 \\
Real-world only & 22 \\
Both real-world and simulation & 71 \\
\end{longtable}
}

The current corpus contains a substantial simulation component: 78
simulation-only records and 71 records reporting both types of
evaluation. This supports a recurring concern about transfer from
controlled simulation to physical flight, especially under illumination
change, occlusion, vibration, communication loss, and degraded sensing.

The structured metric labels occur as follows: ATE/RMSE (170), success
rate (150), computational latency (114), relative drift percentage
(110), and RMSE (45). These labels are multi-label. Although ATE/RMSE is
frequently marked as a reported metric, the numeric
\texttt{ate\_rmse\_m} field is \texttt{NOT\_REPORTED} or otherwise
unavailable for most records. Accordingly, this review does not claim a
pooled meter-level accuracy or rank methods by accuracy.

\subsubsection{3.7 Quality appraisal}\label{37-quality-appraisal}

{\def\LTcaptype{none} % do not increment counter
\begin{longtable}[]{@{}lr@{}}
\toprule\noalign{}
QA tier & Papers \\
\midrule\noalign{}
\endhead
\bottomrule\noalign{}
\endlastfoot
Q-high & 38 \\
Q-medium & 84 \\
Q-low & 49 \\
\textbf{Total} & \textbf{171} \\
\end{longtable}
}

Citation tiers are Core (35), Important (87), and Peripheral (49). The
Core/Important working subset therefore contains 122 records. The QA
distribution indicates that most currently extracted papers report
enough information to receive Q-high or Q-medium status under the
repository rubric, but the scores remain machine-assisted and pending
sample-level human review.

\subsubsection{3.8 Multi-agent systems}\label{38-multi-agent-systems}

The extracted field identifies 85 multi-agent and 86 single-agent
records. Multi-agent status should be interpreted cautiously because
terminology such as swarm, cooperative, or multi-robot can be detected
from text without proving that scalable cooperative localization was
experimentally demonstrated. The key unresolved evidence needs are
communication reliability, synchronization, relative-observation
quality, network scale, and failure recovery.

\subsection{4. Discussion}\label{4-discussion}

\subsubsection{4.1 RQ1: Sensor-fusion
configurations}\label{41-rq1-sensor-fusion-configurations}

The Version-1 corpus is dominated by combinations of inertial sensing
with vision and/or LiDAR. IMU appears in 170 of 171 records, while 3-D
LiDAR and stereo cameras appear in 93 and 87 records, respectively. This
pattern is consistent with complementary observability: inertial data
provide short-term motion propagation, cameras provide visual
constraints, and LiDAR provides range geometry. Radio/UWB, radar, sonar,
optical flow, thermal cameras, and depth cameras occur less often but
represent targeted responses to specific failure modes or operational
constraints.

The frequency pattern answers which configurations are common, but not
which configuration is best. A valid comparative answer requires matched
datasets, common ground truth, comparable trajectories, and consistent
reporting of compute and environmental conditions, which are not
uniformly available in the current table.

\subsubsection{4.2 RQ2: Metrics and
accuracy}\label{42-rq2-metrics-and-accuracy}

The records report a heterogeneous collection of ATE, RMSE, relative
drift, success, and latency indicators. The distribution shows
substantial attention to trajectory and operational metrics, but the
structured numeric field is mostly missing or not reported. Differences
in units, sequences, trajectories, flight speed, map scale, ground-truth
source, and failure definitions prevent direct numerical pooling. The
defensible conclusion is therefore that the field commonly evaluates
trajectory error, drift, success, and computation, not that any one
method achieves a universal accuracy level.

\subsubsection{4.3 RQ3: Algorithmic approaches and comparative
advantages}\label{43-rq3-algorithmic-approaches-and-comparative-advantages}

Visual-LiDAR-inertial fusion is the largest extracted method category,
followed by visual-inertial odometry. LiDAR-inertial and radio-based
approaches form smaller but technically important groups.
Visual-inertial systems offer compactness and high-rate updates;
LiDAR-inertial systems can strengthen geometric observability where
visual texture is poor; radio-based systems can provide external
constraints where infrastructure is available. The comparative advantage
of each approach is conditional on environment, payload, latency,
calibration, and failure mode. The present evidence does not justify
declaring a universal winner.

\subsubsection{4.4 RQ4: Limitations and future
directions}\label{44-rq4-limitations-and-future-directions}

The most important limitations are incomplete full-text availability,
heterogeneous evaluation, limited numeric reporting, and the
simulation-to-real gap. Future studies should report common benchmark
sequences, sensor calibration, synchronization, ground-truth source,
trajectory length, runtime, power or payload constraints, failure cases,
and matched baseline conditions. Field evaluations should include
deliberate degradation, such as lighting change, motion blur, dust,
foliage, occlusion, magnetic or radio disturbance, and communication
loss. Multi-agent studies should report network scale, bandwidth,
latency, packet loss, synchronization, and recovery behavior.

\subsection{5. Limitations and Threats to
Validity}\label{5-limitations-and-threats-to-validity}

First, only 171 of 636 included papers had local PDFs at the time of
extraction. Access and retrieval availability can bias the Version-1
corpus toward recent, open, or institutionally accessible work. Second,
the extraction is machine-assisted. Sensor, environment, method, and
metric labels may contain false positives or omissions; the 20-paper
validation sample must be manually completed before the results are
treated as final coded evidence. Third, the studies use incompatible
datasets, hardware, trajectories, metrics, and definitions of success,
so pooled effect sizes and universal rankings are not appropriate.
Fourth, multi-label fields produce counts greater than the number of
papers and cannot be interpreted as mutually exclusive distributions.
Finally, the 2026 portion of the corpus and the retrieval state are time
sensitive and should be regenerated when the search or full-text set
changes.

\subsection{6. Conclusion}\label{6-conclusion}

This Version-1 full-text synthesis establishes a reproducible baseline
for reviewing GPS/GNSS-denied UAV navigation. The available 171 records
show a strong concentration of visual-inertial and visual-LiDAR-inertial
approaches, frequent use of IMU, camera, and LiDAR sensing, and
substantial reliance on simulation or mixed evaluation. They also show
that the literature commonly reports trajectory, drift, success, and
computational metrics but does not provide a sufficiently uniform
numeric basis for pooled accuracy ranking.

The conclusions are intentionally limited to the currently available
full-text subset. When additional PDFs arrive, the Phase-7 extraction
runner, validation sample, core subset, figures, synthesis matrices, and
manuscript must be regenerated. The final review should incorporate the
enlarged corpus and completed human validation.

\subsection{7. Data and Code
Availability}\label{7-data-and-code-availability}

The analysis tables and scripts are stored in the repository. The
principal files are:

\begin{itemize}
\tightlist
\item
  \texttt{02\_data\_processed/extracted\_master\_v2.csv}
\item
  \texttt{02\_data\_processed/screened\_included\_v2.csv}
\item
  \texttt{02\_data\_processed/screened\_included\_v2\_fulltext.csv}
\item
  \texttt{02\_data\_processed/core\_papers\_v1.csv}
\item
  \texttt{06\_analysis/scripts/16\_execute\_phase7\_extraction.py}
\item
  \texttt{06\_analysis/scripts/06\_generate\_figures\_v2.py}
\item
  \texttt{06\_analysis/scripts/07\_generate\_synthesis.py}
\item
  \texttt{08\_docs/extraction\_validation\_sample\_v1.csv}
\end{itemize}

Local PDFs are not assumed to be redistributable and may be gitignored
or subject to publisher access restrictions. Study-level DOI metadata
should be resolved against the master CSV and checked before submission.

\subsection{8. Figure Captions}\label{8-figure-captions}

\textbf{Figure 1. Publication trends.} Annual distribution of records in
the Version-1 full-text corpus; this is not the trend for all 636
included records.

\textbf{Figure 2. Platform distribution.} Extracted platform labels,
including multi-platform records.

\textbf{Figure 3. Environment distribution.} Multi-label environment
frequencies; categories are not mutually exclusive.

\textbf{Figure 4. Method evolution.} Distribution of extracted primary
methods by publication year in the available full-text corpus.

\textbf{Figure 5. Sensor frequency.} Multi-label sensor frequencies,
showing the co-occurrence of inertial, camera, LiDAR, radio, radar,
sonar, and related modalities.

\textbf{Figure 6. Application domains.} Extracted application-domain
labels.

\textbf{Figure 7. PRISMA-style flow.} Review-record counts from raw
identification through deduplication, inclusion, and Version-1 full-text
extraction.

\textbf{Figure 8. Method-environment heatmap.} Descriptive co-occurrence
of primary method categories and environment labels.

\textbf{Figure 9. Research-maturity radar.} Descriptive quality and
evaluation dimensions for the available full-text corpus; not a causal
or predictive score.

\subsection{9. Conflicts of Interest and
Funding}\label{9-conflicts-of-interest-and-funding}

No conflict-of-interest or funding information was supplied in the
repository materials used for this Version-1 synthesis. These statements
require author confirmation before submission.

\subsection{10. References}\label{10-references}

The following verified methodological and benchmark references support
the review method and representative navigation baselines.
Included-study metadata remain in the authoritative extraction CSV.

\begin{enumerate}
\def\labelenumi{\arabic{enumi}.}
\tightlist
\item
  M. J. Page et al., ``The PRISMA 2020 statement: An updated guideline
  for reporting systematic reviews,'' \emph{BMJ}, vol. 372, p. n71,
  2021, doi: 10.1136/bmj.n71.
\item
  M. L. Rethlefsen et al., ``PRISMA-S: An extension to the PRISMA
  statement for reporting literature searches in systematic reviews,''
  \emph{Syst. Rev.}, vol. 10, no. 1, p. 39, 2021.
\item
  C. Campos et al., ``ORB-SLAM3: An accurate open-source library for
  visual, visual-inertial, and multimap SLAM,'' \emph{IEEE Trans.
  Robot.}, vol. 37, no. 6, pp. 1874--1890, 2021, doi:
  10.1109/TRO.2021.3075644.
\item
  T. Qin, P. Li, and S. Shen, ``VINS-Mono: A robust and versatile
  monocular visual-inertial state estimator,'' \emph{IEEE Trans.
  Robot.}, vol. 34, no. 4, pp. 1004--1020, 2018, doi:
  10.1109/TRO.2018.2853729.
\item
  J. Zhang and S. Singh, ``LOAM: Lidar odometry and mapping in
  real-time,'' \emph{Robotics: Science and Systems}, 2014, doi:
  10.15607/RSS.2014.X.007.
\item
  W. Xu et al., ``FAST-LIO2: Fast direct lidar-inertial odometry,''
  \emph{IEEE Trans. Robot.}, vol. 38, no. 4, pp. 2053--2073, 2022, doi:
  10.1109/TRO.2022.3141876.
\item
  T. Shan et al., ``LIO-SAM: Tightly-coupled lidar inertial odometry via
  smoothing and mapping,'' in \emph{Proc. IEEE/RSJ IROS}, 2020, pp.
  5135--5142, doi: 10.1109/IROS45743.2020.9341176.
\item
  M. Burri et al., ``The EuRoC micro aerial vehicle datasets,''
  \emph{Int. J. Robot. Res.}, vol. 35, no. 10, pp. 1157--1163, 2016,
  doi: 10.1177/0278364916652421.
\end{enumerate}

\subsection{Appendix A. Reproducibility and Quality-Control
Audit}\label{appendix-a-reproducibility-and-quality-control-audit}

{\def\LTcaptype{none} % do not increment counter
\begin{tabularx}{\columnwidth}{@{}Xl@{}}
\toprule\noalign{}
Check & Result \\
\midrule\noalign{}
\endhead
\bottomrule\noalign{}
\endlastfoot
\texttt{extracted\_master\_v2.csv} row count = 171 & PASS \\
\texttt{screened\_included\_v2.csv} row count = 636 & PASS \\
Full-text subset row count = 171 & PASS \\
Full-text coverage = 171/636 = 26.9\% & PASS \\
QA totals within 0--10 & PASS \\
QA tiers valid & PASS \\
QA counts = 38 / 84 / 49 & PASS \\
Experiment counts = 22 / 78 / 71 & PASS \\
Citation tiers = 35 / 87 / 49 & PASS \\
Core + Important = 122 & PASS \\
Human validation sample completed & REVIEW REQUIRED \\
Complete final-corpus claim justified & NO \\
\end{tabularx}
}

\textbf{Overall status: REVIEW REQUIRED.} The numerical Version-1 checks
pass, but human validation and additional full-text retrieval remain
incomplete. The manuscript is suitable as an intermediate synthesis and
reproducibility baseline, not as the final complete-corpus submission.

\end{document}
